\documentclass[11pt,a4paper]{article}
\usepackage[margin=2.5cm]{geometry}
\usepackage[T1]{fontenc}
\usepackage[utf8]{inputenc}
\usepackage{lmodern}
\usepackage{microtype}
\usepackage[hidelinks]{hyperref}
\setlength{\parindent}{0pt}
\setlength{\parskip}{0.7em}

\begin{document}

\begin{flushright}
October 2026
\end{flushright}

\noindent
ChemRxiv Editorial Team

\vspace{1.2em}
\noindent
\textbf{Re: Submission of ``Beyond Candidate Generation: A Validation-First Protocol for AI-Assisted Transition-State Discovery''}

\vspace{0.6em}
Dear Editors,

I am pleased to submit the manuscript above for consideration as a ChemRxiv
preprint. The manuscript is a methods and perspective contribution on how to
validate transition-state (TS) candidates produced by AI models. It is not a
claim of a new predictive model; it is a protocol that separates candidate
generation from physical validation and makes the validation layer
reproducible.

The central argument is that a generated geometry is a hypothesis, not a
transition state. A TS requires one imaginary frequency, the intended atom
pair, the correct connectivity, and a consistent electronic state. The
manuscript presents a five-stage protocol---chemical preflight, barrier-region
scan, constrained reference plus Hessian, OptTS, and validation---and implements it as
a portable, open-source skill package for local ORCA/xtb and Gaussian/Slurm
execution. It illustrates the protocol with internal case studies from Ni/Fe
catalysis, including a per-point scan strategy that completed 16/16 points
after a collective scan failed three times, and a generative-model test on an
84-atom broken-symmetry organometallic system where the candidates were out of
distribution.

I believe the work is well suited to ChemRxiv because it is a community-facing
methods contribution: the protocol and the implementation are meant to be used
and criticized before journal submission. The preprint will also establish a
citable reference for the open-source package. The manuscript is not under
consideration elsewhere, and I am the sole author.

The author declares no competing financial interest. The code and manuscript source are available at \url{https://github.com/lhfjjk/ts-computation-suite}.

Suggested categories: Computational Chemistry (primary); Catalysis (secondary).

\vspace{1.0em}
\noindent
Sincerely,\\[0.8em]
Hao-Fan Li\\
State Key Laboratory of Structural Chemistry\\
Fujian Institute of Research on the Structure of Matter\\
Chinese Academy of Sciences, Fuzhou 350002, China\\
University of Chinese Academy of Sciences, Beijing 100049, China\\
\texttt{[corresponding email]}

\end{document}
